\newcommand{\BrownianCircuitCaption}{%
Forcing exit near an arbitrary boundary. (a,b) The first path segment crosses the rectangle $Q$ from bottom to top. It then follows an ordered, sufficiently thin corridor for one counterclockwise turn around $\overline B(0,2)$ and returns across $Q$ from left to right. The two crossings intersect, creating a closed subpath with winding number one about the disk; its trace separates the disk from infinity. Self-intersections are allowed. Thus the path meets every connected unbounded complement that meets $\overline B(0,1)$. (c) At each entrance $\sigma_j$ after the preceding block, a block of length $Ar^2$ begins. The orange pieces represent occupation of $E_r=\{u\in D:\operatorname{dist}(u,\partial D)\le r\}$; the gaps have none. Conditional on each block start before exit $\zeta$, the probability of exit in that block is at least $\eta$. Hence the expected occupation is at most $Ar^2/\eta$.%
}

\newcommand{\RibbonFrontCaption}{%
One step in ribbon coordinates $P(s,x)$. The black past stays in the core $|x|<\epsilon$ and first reaches the front $s=k$ at $a$, with tip offset at most $\epsilon/4$. The blue access guide moves to the centerline and passes through the interior waypoint $P(k+1/2,0)$. The orange region is $P((k+1/3,k+2/3)\times(-\epsilon,\epsilon))$, where the next step's sausage area is tested before its first crossing of $s=k+1$. The guide depicts an available route, not the random continuation. The transverse scale and core width are exaggerated.%
}

\newcommand{\SerpentineRibbonCaption}{%
A dense visit through a narrow ribbon. (a) In normalized half-plane coordinates, the marked point $w$ is separated from a disk of radius $s_0/20$ centered at the horizontally displaced point $w'$. The connector starts vertically at $0$, reaches the first strand from below, and stays away from $w$; the final strand continues upward. (b) The disk contains $2m+1$ strands of length $2h_0$ joined by alternating exterior semicircles. Choosing ribbon half-width $\ell=h_0/(C_0m)$ gives order $m^2$ steps. The orange marks indicate selected disjoint middle test regions. Their sausage tests contribute order $m^2\ell^d\asymp m^{2-d}=m^p$. Both panels are schematic and not to scale; the array in (b) is enlarged, and the ribbon widths are exaggerated.%
}
